\section{Introduction}

The application of Neutrosophic Logic to LLM uncertainty quantification
(Leyva-Vázquez \& Smarandache, 2025) represents an important departure
from probabilistic frameworks. By treating Truth (T), Indeterminacy (I),
and Falsity (F) as independent dimensions, the neutrosophic approach
allows models to express internal conflicts that probabilistic
constraints (T+I+F=1) suppress. Their finding that 35\% of evaluations
exhibit ``hyper-truth'' (T+I+F > 1.0) under unconstrained
prompting demonstrates that LLMs carry richer epistemic information than
probabilistic output formats can express.

However, their study has three limitations that constrain the
generalizability of their findings:

\begin{enumerate}
\def\labelenumi{\arabic{enumi}.}
\item
  \textbf{Single vendor}: All four models (GPT-4o, GPT-4-turbo,
  GPT-3.5-turbo, GPT-4o-mini) come from OpenAI. The observed hyper-truth
  could be an artifact of OpenAI's training pipeline rather than a
  general property of LLMs.
\item
  \textbf{Single shot}: Each cell (model × phenomenon × strategy) was
  evaluated once, providing no measure of intra-model consistency or
  statistical significance.
\item
  \textbf{Unpublished prompts}: The specific prompts used for each
  strategy were not committed to the repository, preventing independent
  replication.
\end{enumerate}

We address all three limitations. More importantly, we identify a fourth
problem that scalar T/I/F cannot solve even in principle: \textbf{the
Absorption problem}, where some models collapse all uncertain states to
(T=0, I=1, F=0), losing the distinction between types of uncertainty
that neutrosophic logic was designed to preserve.

We then demonstrate that this problem is solved by extending the output
format from scalars to tensors: specifically, by requiring models to
declare structured losses alongside their T/I/F values.
